% Finalizado 0 
\chapter{Justificación y Objetivos}
\label{sec:justificacion_y_objetivos}

\section{Motivación}
Efectos como el fuego, el agua o el humo son indispensables para crear escenas 3D realistas. Por desgracia, esto consume muchos recursos debido a la gran cantidad de datos.
Las APIs más antiguas saturan la CPU con tareas de control (\textit{CPU overhead}). Esto frena el rendimiento general del sistema.

La meta de este proyecto es renderizar nubes de partículas  masivas. Para ello usaremos una API moderna llamada Vulkan. 
Vulkan nos permite comunicarnos directamente con el hardware sin intermediarios pesados. Esto nos abre las puertas a densidades de elementos muy altas. 
Para ello necesitamos  aplicar leyes físicas en geometría simple para obtener resultados rápidos. Al final debemos lograr un equilibrio entre resultado y rendimiento.

 
\section{Objetivos}
El objetivo consiste en crear un motor de partículas rápido y eficiente con Vulkan. Para ello, nos marcamos las siguientes tareas concretas:                                     

\begin{itemize}
    \item \textbf{Cómputo en GPU:} La tarjeta gráfica asume todo el trabajo real. Programamos \textit{Compute Shaders} para que calculen el ciclo de vida, la posición y el color de cada punto. La CPU ya no interviene aquí.

    \item \textbf{Fluidos y gases:} Escribir algoritmos de tipo procedimental directos en el hardware gráfico. Con esto imitamos la física real detrás de líquidos o columnas de humo.                               

    \item \textbf{Efectos con volumen:} Dar sensación de masa y densidad al fuego. Lo haremos sumando técnicas de mezcla de color (\textit{Alpha Blending}) con el uso de \textit{Point Sprites}.

    \item \textbf{Sincronización:} Configurar barreras de memoria y organizar bien las colas en Vulkan. Esto evita condiciones de carrera (\textit{race conditions}) al procesar la información en paralelo, impidiendo colisiones entre operaciones de lectura y escritura simultáneas.

    \item \textbf{Control estético:} Crear materiales sencillos que cambien con el tiempo. El brillo, el color y la opacidad se alteran solos según la edad de la partícula.
\end{itemize}